\documentclass[11pt]{article}
\usepackage[T1]{fontenc}
\usepackage[utf8]{inputenc}
\usepackage{lmodern}
\usepackage[margin=1in]{geometry}
\usepackage{amsmath,amssymb}
\usepackage{booktabs}
\usepackage[round,authoryear]{natbib}
\usepackage[hidelinks]{hyperref}
\usepackage{xcolor}
\usepackage{setspace}
\setstretch{1.08}
% generated by `python -m voc.note` from results/; do not edit
\newcommand{\nAInBH}{0.53}
\newcommand{\nAInBHSE}{0.11}
\newcommand{\nAInEnsSE}{0.12}
\newcommand{\nAInEnsSR}{0.31}
\newcommand{\nAInEnsT}{1.1}
\newcommand{\nAInIR}{0.11}
\newcommand{\nAInLinSR}{0.39}
\newcommand{\nAInLinT}{1.9}
\newcommand{\nAInLong}{719}
\newcommand{\nAInN}{1080}
\newcommand{\nAInSR}{0.31}
\newcommand{\nAInT}{1.1}
\newcommand{\nAPostBH}{0.57}
\newcommand{\nAPostBHSE}{0.47}
\newcommand{\nAPostEnsSE}{0.45}
\newcommand{\nAPostEnsSR}{-0.01}
\newcommand{\nAPostEnsT}{0.1}
\newcommand{\nAPostIR}{0.06}
\newcommand{\nAPostLinSR}{0.30}
\newcommand{\nAPostLinT}{0.1}
\newcommand{\nAPostLong}{34}
\newcommand{\nAPostN}{60}
\newcommand{\nAPostSR}{-0.01}
\newcommand{\nAPostT}{0.1}
\newcommand{\nBDraws}{300}
\newcommand{\nBInBH}{0.51}
\newcommand{\nBInBHSE}{0.11}
\newcommand{\nBInEnsSE}{0.09}
\newcommand{\nBInEnsSR}{0.50}
\newcommand{\nBInEnsT}{3.1}
\newcommand{\nBInIR}{0.33}
\newcommand{\nBInLinSR}{0.46}
\newcommand{\nBInLinT}{3.1}
\newcommand{\nBInLong}{720}
\newcommand{\nBInN}{1080}
\newcommand{\nBInSR}{0.49}
\newcommand{\nBInT}{3.1}
\newcommand{\nBPostBH}{0.71}
\newcommand{\nBPostBHSE}{0.47}
\newcommand{\nBPostEnsSE}{0.45}
\newcommand{\nBPostEnsSR}{-0.06}
\newcommand{\nBPostEnsT}{-0.4}
\newcommand{\nBPostIR}{-0.17}
\newcommand{\nBPostLinSR}{0.45}
\newcommand{\nBPostLinT}{0.1}
\newcommand{\nBPostLong}{39}
\newcommand{\nBPostN}{60}
\newcommand{\nBPostSR}{-0.06}
\newcommand{\nBPostT}{-0.4}
\newcommand{\nCorrDfy}{0.76}
\newcommand{\nCorrDp}{0.76}
\newcommand{\nCorrEp}{0.82}
\newcommand{\nCorrInfl}{0.52}
\newcommand{\nCorrLty}{0.65}
\newcommand{\nCorrR}{0.98}
\newcommand{\nCorrRet}{0.992}
\newcommand{\nCtry}{41}
\newcommand{\nCtryVm}{31}
\newcommand{\nCtryWins}{7}
\newcommand{\nDecABH}{0.53}
\newcommand{\nDecAIR}{0.11}
\newcommand{\nDecALinSR}{0.39}
\newcommand{\nDecALinT}{1.9}
\newcommand{\nDecASR}{0.31}
\newcommand{\nDecAT}{1.0}
\newcommand{\nDecDraws}{10}
\newcommand{\nDecFourBH}{0.51}
\newcommand{\nDecFourIR}{0.33}
\newcommand{\nDecFourLinSR}{0.46}
\newcommand{\nDecFourLinT}{3.1}
\newcommand{\nDecFourSR}{0.49}
\newcommand{\nDecFourT}{3.1}
\newcommand{\nDecOneBH}{0.53}
\newcommand{\nDecOneIR}{0.14}
\newcommand{\nDecOneLinSR}{0.39}
\newcommand{\nDecOneLinT}{1.9}
\newcommand{\nDecOneSR}{0.34}
\newcommand{\nDecOneT}{1.3}
\newcommand{\nDecStepMax}{0.07}
\newcommand{\nDecStepMin}{0.03}
\newcommand{\nDecThreeBH}{0.52}
\newcommand{\nDecThreeIR}{0.23}
\newcommand{\nDecThreeLinSR}{0.42}
\newcommand{\nDecThreeLinT}{2.3}
\newcommand{\nDecThreeSR}{0.42}
\newcommand{\nDecThreeT}{2.1}
\newcommand{\nDecTwoBH}{0.53}
\newcommand{\nDecTwoIR}{0.18}
\newcommand{\nDecTwoLinSR}{0.39}
\newcommand{\nDecTwoLinT}{1.9}
\newcommand{\nDecTwoSR}{0.38}
\newcommand{\nDecTwoT}{1.6}
\newcommand{\nExtCommit}{9c77321}
\newcommand{\nExtDraws}{10}
\newcommand{\nExtIR}{0.33}
\newcommand{\nExtLinIR}{0.33}
\newcommand{\nExtLinSR}{0.46}
\newcommand{\nExtLinT}{3.1}
\newcommand{\nExtLinZero}{-0.11}
\newcommand{\nExtSR}{0.49}
\newcommand{\nExtSeedSD}{0.013}
\newcommand{\nExtT}{3.1}
\newcommand{\nExtraLagSR}{0.29}
\newcommand{\nHybPredSR}{0.40}
\newcommand{\nHybRetSR}{0.36}
\newcommand{\nKmzIR}{0.31}
\newcommand{\nKmzLinIR}{0.33}
\newcommand{\nKmzLinSR}{0.46}
\newcommand{\nKmzLinT}{3.1}
\newcommand{\nKmzLinZero}{-0.11}
\newcommand{\nKmzSR}{0.47}
\newcommand{\nKmzT}{2.9}
\newcommand{\nLaBH}{0.66}
\newcommand{\nLaIR}{0.21}
\newcommand{\nLaLinSR}{0.60}
\newcommand{\nLaLinT}{3.1}
\newcommand{\nLaLinZero}{-0.11}
\newcommand{\nLaSR}{0.46}
\newcommand{\nLaT}{1.9}
\newcommand{\nLagBH}{0.53}
\newcommand{\nNagelRawIR}{0.26}
\newcommand{\nNagelRawT}{2.42}
\newcommand{\nNkCommit}{40ab1c8}
\newcommand{\nNkDraws}{500}
\newcommand{\nNkSR}{0.42}
\newcommand{\nOffBH}{0.58}
\newcommand{\nOffN}{61}
\newcommand{\nOffSR}{0.00}
\newcommand{\nOffT}{0.1}
\newcommand{\nOffVmSR}{0.04}
\newcommand{\nOwnIR}{0.11}
\newcommand{\nOwnLinIR}{0.20}
\newcommand{\nOwnLinSR}{0.39}
\newcommand{\nOwnLinT}{1.9}
\newcommand{\nOwnLinZero}{-0.16}
\newcommand{\nOwnSR}{0.31}
\newcommand{\nOwnT}{1.0}
\newcommand{\nPoolT}{0.2}
\newcommand{\nPostFirst}{2021-01}
\newcommand{\nPostLast}{2025-12}
\newcommand{\nRawIR}{0.29}
\newcommand{\nRawT}{2.8}
\newcommand{\nRawZeroIR}{0.27}
\newcommand{\nRawZeroT}{2.57}
\newcommand{\nRepDraws}{300}
\newcommand{\nRepIR}{0.11}
\newcommand{\nRepSR}{0.31}
\newcommand{\nRepSeedSD}{0.017}
\newcommand{\nRepT}{1.1}
\newcommand{\nRmsSR}{0.35}
\newcommand{\nRrSR}{0.495}

\hypersetup{pdftitle={Four Data Choices and the Virtue of Complexity}, pdfauthor={Mohy Mabrouk},
  pdfkeywords={return predictability, random features, machine learning, market timing, replication}}
\newcommand{\Rplus}{R_{t+1}}

\title{Four Data Choices and the Virtue of Complexity:\\ A Replication Note on Kelly, Malamud and Zhou (2024)}
\author{Mohy Mabrouk\thanks{Code, data manifests and every number in this note:
\url{https://github.com/mohymabrouk/virtue-of-complexity-out-of-sample}. The code and a first draft of the text were
prepared with the help of an AI coding assistant (Claude, Anthropic); the author checked every number against the
repository's result files and is responsible for all content and errors.}}
\date{November 2026}

\begin{document}
\maketitle

\begin{abstract}
\noindent Kelly, Malamud and Zhou (2024) report an out-of-sample Sharpe ratio of \nKmzSR{} for market timing
with 12{,}000 random Fourier features of the Welch--Goyal predictors, and write that their findings are insensitive
to how standardisations are implemented. Holding the estimation engine fixed, I show that four data-construction
choices the paper does not pin down move this Sharpe ratio from \nAInSR{} to \nBInSR{} and its alpha t-statistic
from \nAInT{} to \nBInT: where the predictors' expanding standardisation starts, the inflation publication lag, the
market-return series, and centred versus uncentred return volatility. One choice uses inflation before its
release. A volatility window that includes the forecast month, used in one public replication, matches the
published Sharpe ratio but not the alpha. Scored once over 2021--2025, the construction that reproduces the paper
earns a Sharpe ratio of $\nBPostEnsSR$ against \nBPostBH{} for buy-and-hold.

\medskip
\noindent\textbf{JEL:} G11, G12, G17, C53.\quad \textbf{Keywords:} return predictability, random features,
machine learning, market timing, replication.
\end{abstract}

\section{Introduction}

\citet{kmz2024} (KMZ) time the US stock market with ridge regressions on up to 12{,}000 random Fourier features
(RFF) of the 15 Welch--Goyal predictors, estimated on rolling 12-month windows. Their headline strategy
($P = 12{,}000$, ridge penalty $z = 10^3$) has an out-of-sample Sharpe ratio of \nKmzSR{} over 1931--2020 and an
alpha t-statistic of \nKmzT{} against the volatility-standardised market (their Table~1). On data construction,
the paper says:
\begin{quote}
\small ``We volatility standardize returns and predictors using backward-looking standard deviations that preserve
the out-of-sample nature of our forecasts. Returns are standardized by their trailing 12-month return standard
deviation [\dots]. In contrast, predictors are standardized using an expanding window historical standard deviation
[\dots]. Still, our results are insensitive to this step (none of our findings are sensitive to variations in how
standardizations are implemented).'' (KMZ, Section 6.1)
\end{quote}
KMZ's code is not public, and public replications disagree on the level of the headline Sharpe ratio. A
replication that follows the text above with a lagged 12-month standard deviation \citep{thisrepo} obtains
\nRepSR{} (\nRepDraws{} draws; alpha t \nRepT); one public replication \citep{nkonts} obtains \nNkSR{} with a
volatility window that includes the month being predicted; another \citep{rrikrv} reports \nRrSR.

This note holds the estimation engine fixed and explains the spread. Three findings:
\begin{enumerate}\itemsep0pt
\item Four data-construction choices, none of which the paper pins down, move the Sharpe ratio from \nDecASR{} to
\nDecFourSR{} and the alpha t-statistic from \nDecAT{} to \nDecFourT{} (Table~\ref{tab:decomp}). With all four the
engine reproduces KMZ's nonlinear \emph{and} linear rows. Three of the choices use only past data; the fourth uses
inflation without its usual one-month publication lag.
\item A volatility window that includes the forecast month gives a \emph{partial} match (Sharpe \nLaSR{}, alpha t
\nLaT{}). Matching the Sharpe ratio alone is therefore weak evidence that a replication is right.
\item Scored once after the paper's sample (\nPostFirst{} to \nPostLast), the construction that reproduces the
paper earns a Sharpe ratio of $\nBPostEnsSR$ against \nBPostBH{} for buy-and-hold; the pre-registered one earns
$\nAPostEnsSR$ against \nAPostBH. Reproducing the in-sample number does not rescue the out-of-sample one.
\end{enumerate}
I make no claim about which choices KMZ's own code makes. The point is narrower: the published level is reproduced
by one reasonable construction and missed by others that are at least as close to the paper's text, so the
headline number is more sensitive to data construction than the quoted sentence suggests.

\section{Setup}

The engine follows KMZ and is the same in every row below. Each month $t$ the 15 scaled predictors $G_t$ are
mapped to $P/2$ pairs $\sin(\gamma\,\omega_i^{\top} G_t), \cos(\gamma\,\omega_i^{\top} G_t)$ with $\omega_i \sim
N(0, I)$ and $\gamma = 2$. Each feature is divided by its standard deviation over the 12 training months; a ridge
regression without intercept, penalty $z T$ in dual form, is fitted on the pairs $(S_{s}, R_{s+1})$,
$s = t-12, \dots, t-1$, and the forecast for $R_{t+1}$ is the position. Statistics are computed per draw of
$\omega$ and averaged over draws, as in KMZ; the alpha is from an OLS regression of the strategy on the scaled
market return. Targets run from 1931-01 to 2020-12 (1{,}080 months). The dual engine equals a direct primal ridge
to $10^{-8}$, and on raw (unscaled) returns its ridgeless strategy has an alpha t-statistic of \nRawZeroT{}
and an information ratio of \nRawZeroIR, against \nNagelRawT{} and \nNagelRawIR{} that \citet{nagel2025} obtains
with KMZ's own code. Section 3 shows the rest of the gap is in the data, not in the machine.

\section{Four data choices}

Construction A is the pre-registered one from \citet{thisrepo}. Construction B changes four steps, one at a
time:
\begin{enumerate}\itemsep0pt
\item[(i)] \textbf{Start of the predictors' expanding standardisation.} A starts it in 1926-07, the first CRSP
month; B at each series' first month in the Welch--Goyal file (1871 for the valuation ratios). The scaled series
differ a lot in the early sample: on common months the correlation between the two versions is \nCorrDp{} for
the dividend-price ratio and \nCorrLty{} for the long-term yield. Because predictors enter the features through
$\gamma\,\omega^{\top} G_t$, their scale acts like a kernel bandwidth.
\item[(ii)] \textbf{Inflation lag.} CPI for month $t$ is published during month $t+1$. \citet{welchgoyal2008}
therefore lag it one month, and A does too; B uses it unlagged. This is the one choice that uses information not
available at the end of month $t$ (by about two weeks). KMZ do not say which they use.
\item[(iii)] \textbf{Market return.} A uses Ken French's Mkt-RF (CRSP value-weighted, all stocks); B uses the
S\&P~500 return with dividends minus the risk-free rate from the Welch--Goyal file (correlation \nCorrRet).
\item[(iv)] \textbf{Return volatility.} A divides $r_{t}$ by the standard deviation of $r_{t-12},\dots,r_{t-1}$, as
the text says; B by the uncentred root mean square of the same twelve returns, which is what \citet{guo2026} and
\citet{rrikrv} use.
\end{enumerate}
Table~\ref{tab:decomp} adds the switches cumulatively. Each adds between \nDecStepMin{} and \nDecStepMax{} to the Sharpe ratio; none
closes the gap alone. With all four, the engine gives a Sharpe ratio of \nDecFourSR{} (alpha t \nDecFourT, IR \nDecFourIR)
against KMZ's \nKmzSR{} (\nKmzT, \nKmzIR), and the linear kitchen-sink row matches too (\nDecFourLinSR{} and
\nDecFourLinT{} against \nKmzLinSR{} and \nKmzLinT). Running the same engine on the data file committed in
\citet{rrikrv} gives the same numbers (\nExtSR, \nExtT, \nExtIR), so the four switches account for the whole
difference between the two data files. Swapping only that file's returns into A gives \nHybRetSR; swapping only its
predictors gives \nHybPredSR: both blocks matter.

\begin{table}[t]
\centering\small
\caption{Same engine, different data construction (US, targets 1931-01 to 2020-12, $P = 12{,}000$,
$z = 10^3$, \nDecDraws{} draws per row, statistics averaged over draws; with 300 draws A gives \nAInSR{} and B
\nBInSR, Table~\ref{tab:post}). B\&H = the scaled market.}
\label{tab:decomp}
\begin{tabular}{lcccccc}
\toprule
 & \multicolumn{3}{c}{RFF strategy} & \multicolumn{2}{c}{Linear, $z=10^3$} & B\&H \\
\cmidrule(lr){2-4}\cmidrule(lr){5-6}
Construction & Sharpe & $\alpha$ t & IR & Sharpe & $\alpha$ t & Sharpe \\
\midrule
A (pre-registered) & \nDecASR & \nDecAT & \nDecAIR & \nDecALinSR & \nDecALinT & \nDecABH \\
+ (i) predictor std from first month in file & \nDecOneSR & \nDecOneT & \nDecOneIR & \nDecOneLinSR & \nDecOneLinT & \nDecOneBH \\
+ (ii) inflation not lagged & \nDecTwoSR & \nDecTwoT & \nDecTwoIR & \nDecTwoLinSR & \nDecTwoLinT & \nDecTwoBH \\
+ (iii) S\&P 500 return (Welch--Goyal file) & \nDecThreeSR & \nDecThreeT & \nDecThreeIR & \nDecThreeLinSR & \nDecThreeLinT & \nDecThreeBH \\
+ (iv) uncentred RMS volatility $=$ B & \nDecFourSR & \nDecFourT & \nDecFourIR & \nDecFourLinSR & \nDecFourLinT & \nDecFourBH \\
\addlinespace
\citet{rrikrv} data file, same engine & \nExtSR & \nExtT & \nExtIR & \nExtLinSR & \nExtLinT & -- \\
KMZ, Table 1 & \nKmzSR & \nKmzT & \nKmzIR & \nKmzLinSR & \nKmzLinT & -- \\
\midrule
A with look-ahead window (3 draws) & \nLaSR & \nLaT & \nLaIR & \nLaLinSR & \nLaLinT & \nLaBH \\
\bottomrule
\end{tabular}
\end{table}

\section{A look-ahead window gives a partial match}

\citet{nkonts} scales the return to be predicted by \texttt{returns.rolling(12).std()}, a window that ends with,
and so contains, $\Rplus$. Applied to construction A, this gives a Sharpe ratio of \nLaSR, close to KMZ's
\nKmzSR, and the linear ridgeless Sharpe ratio of $\nLaLinZero$ that KMZ report. But the alpha t-statistic is only
\nLaT{} (KMZ: \nKmzT) and the linear $z = 10^3$ Sharpe ratio is \nLaLinSR{} (KMZ: \nKmzLinSR). The mechanism is
simple. When a month's own return is in its scale, an extreme month dominates the window: with the other eleven
returns fixed and $r_{t+1} \to \pm\infty$, the scaled return tends to $\pm\sqrt{12} \approx \pm 3.46$ whatever
its size. Crashes are capped, so the scaled market itself looks better (buy-and-hold Sharpe \nLaBH{} against
\nLagBH{} with the lagged window), and a mostly-long strategy inherits that gain while its alpha against the
equally inflated benchmark rises much less. An earlier version of the replication repository took this partial
match as the likely source of the gap; Table~\ref{tab:decomp} shows it is not needed. The practical lesson for
replicators is to match the Sharpe ratio and the alpha jointly.

This timing issue is different from the one in \citet[Appendix A]{nagel2025}, who shows that dividing returns
by a properly lagged volatility can by itself create spurious predictability, and for that reason works with
unscaled returns.

\section{After the paper's sample}

Construction B was defined after the pre-registered runs, so its only clean test is after the paper's
sample. It was scored once, as pre-declared, on the \nBPostN{} months from \nPostFirst{} to \nPostLast{} (the
Welch--Goyal file ends in 2025-12), with the same 300 draws as construction A; Table~\ref{tab:post} reports the
strategy that averages the forecasts over draws, which is the one an investor could trade (per-draw averages are
within 0.01). In sample, B is within KMZ's numbers; after the sample, both constructions lose to buy-and-hold by
more than half a Sharpe ratio, and B's alpha is negative. The linear ridge does better than the RFF strategy in both.

\begin{table}[t]
\centering\small
\caption{In sample and after the sample ($P = 12{,}000$, $z = 10^3$, \nBDraws{} draws; each construction against
its own scaled market). In sample: statistics averaged over draws, as KMZ. After the sample: the strategy that
averages forecasts over draws, Sharpe ratio with its standard error \citep{lo2002}.}
\label{tab:post}
\begin{tabular}{lcccccc}
\toprule
 & \multicolumn{2}{c}{1931--2020} & \multicolumn{4}{c}{\nPostFirst{} to \nPostLast{} (\nBPostN{} months)} \\
\cmidrule(lr){2-3}\cmidrule(lr){4-7}
Construction & Sharpe & $\alpha$ t & Sharpe [SE] & $\alpha$ t & B\&H & Linear \\
\midrule
A (pre-registered) & \nAInSR & \nAInT & $\nAPostEnsSR$ [\nAPostEnsSE] & \nAPostEnsT & \nAPostBH & \nAPostLinSR \\
B (reproduces KMZ) & \nBInSR & \nBInT & $\nBPostEnsSR$ [\nBPostEnsSE] & $\nBPostEnsT$ & \nBPostBH & \nBPostLinSR \\
KMZ, Table 1 & \nKmzSR & \nKmzT & -- & -- & -- & -- \\
\bottomrule
\end{tabular}
\end{table}

Five years is a short window. The standard error of an annualised Sharpe ratio over \nBPostN{} months is about
\nBPostEnsSE, and the chance of an alpha t-statistic above 1.645 if the true information ratio were KMZ's \nKmzIR{}
is only about one in six. The evidence is that the strategy did not behave as advertised, not that it is
significantly negative. The cross-section points the same way: with construction A's recipe on Jensen--Kelly--Pedersen
predictors \citep{jkp2023}, the strategy beats buy-and-hold in \nCtryWins{} of \nCtry{} non-US markets, Nagel's
volatility-timed momentum beats it in \nCtryVm, and the pooled cross-country alpha t-statistic is \nPoolT{}
\citep{thisrepo}. Those runs use A's centred volatility; whether choice (iv) changes them is untested.

\section{Related work}

\citet{nagel2025} shows that with $T = 12$ the RFF forecasts are close to a volatility-timed momentum rule.
\citet{buncic2025} attributes the virtue of complexity to the zero-intercept restriction and to averaging
performance measures across random draws instead of averaging forecasts. \citet{cartea2025} show that measurement
error in the predictors caps the benefit of complexity, and \citet{carteashi2025} that random-feature gains can be
mechanical. \citet{fallahgoul2025} shows that within-window standardisation changes the kernel the RFF approximate,
and derives sample-size bounds that 12-month windows do not meet. \citet{guo2026} find that ridge and ridgeless
RFF models do not beat the historical mean. This note is complementary: before asking why the strategy works, it
asks how much of the published level depends on data choices the paper does not report.

\section{Conclusion}

With the estimation engine held fixed, four data choices move the headline Sharpe ratio of KMZ from \nDecASR{} to
\nDecFourSR. One of them uses inflation two weeks before it is published. A look-ahead volatility window matches
the Sharpe ratio but not the alpha. The construction that reproduces the paper's in-sample numbers does not
beat buy-and-hold after the sample, and neither does the pre-registered one. Replications of high-dimensional return predictors
should report their data construction as precisely as their model, and match several statistics, not one.

{\sloppy\paragraph{Data and code.} Everything is in the repository cited in the title footnote: the pre-registration
and the dated change log (\texttt{PLAN.md}), the scripts that produce every table
(\texttt{scripts/construction\_b.py}, \texttt{scripts/diag\_external\_data.py},
\texttt{scripts/replication\_diagnostics.py}) and the result files the numbers are read from
(\texttt{paper/numbers.json} is generated from them by \texttt{python -m voc.note}). Data: Welch--Goyal monthly
file (2025 vintage), Ken French data library, JKP global factors. The \citet{rrikrv} data file is read, not
redistributed; its SHA-256 is recorded in \texttt{results/external\_data\_diagnostic.json}.\par}

\bibliographystyle{plainnat}
\bibliography{refs}
\end{document}
